\chapter{Modules and Packages}
\label{ch:modules-packages}

A program that outgrows a single file needs structure. \textbf{Modules} and
\textbf{packages} are how Salam splits a large program into pieces, controls what
each piece exposes to the others, and pulls in the standard library. You have been
quietly using this system every time you wrote \code{import "str"}; now we look at
it properly.

\section{Importing the standard library}

The standard library is organised into modules, each a self-contained bundle of
related functions: \code{math} for numbers, \code{str} for strings, \code{io} for
input and output, and many more. To use one, \code{import} it by name and call its
functions with the module name as a prefix:

\begin{salam}
import "math"

func main:
    println math.Sqrt(144.0)
    println math.MaxI(7, 3)
end
\end{salam}

\begin{progout}
12
7
\end{progout}

The prefix \code{math.} is called \textbf{qualification}, and it does two
useful things. It tells the reader exactly where \code{Sqrt} comes from, and it
keeps names from colliding: \code{math.Max} and your own \code{max} can coexist
without confusion.

\begin{note}
A handful of operations need no import at all: \code{print}, \code{println}, and
\code{len} are built into the language and are always available. Everything else
in the standard library lives in a module you import.
\end{note}

\section{Importing several modules}

You can write one \code{import} per module, or list several module paths after a
single \code{import} no commas, no parentheses, just the paths side by side:

\begin{salam}
import "math" "str" "io"

func main:
    io.Println(str.Upper("salam"))
    println math.Gcd(12, 18)
end
\end{salam}

\begin{progout}
SALAM
6
\end{progout}

\section{Aliasing an import}

Sometimes a module's name is long, or clashes with a name you want to use
yourself. You can give an import a shorter \textbf{alias} by writing the alias
before the module name:

\begin{salam}
import strutil "str" "io"

func main:
    io.Println(strutil.Upper("salam"))
end
\end{salam}

\begin{progout}
SALAM
\end{progout}

Now the \code{str} module is reached as \code{strutil}. Aliases are a matter of
convenience and clarity; use them when they make the code read better, and leave
them off otherwise.

\section{How imports are resolved}

When you write \code{import "math"}, Salam looks for the module \code{math} in the
standard library specifically the file \code{std/math/math.salam}. A nested
module path works the same way: \code{import "db/sqlite"} resolves to
\code{std/db/sqlite/sqlite.salam}. An import whose name ends in \code{.salam},
such as \code{import "helpers.salam"}, is instead treated as a \emph{file relative
to the importing file} which is how you bring in your \emph{own} modules, the
subject of the next section.

\section{Writing your own modules}

A module is just a \code{.salam} file with definitions in it. To split your
program across files, put related functions and types in their own file and import
it where you need them. Two pieces of machinery make this work: the \code{package}
declaration and the \code{pub} keyword.

\subsection{Declaring a package}

Every file belongs to a \textbf{package}, named at the top with the \code{package}
keyword. Your program's entry point lives in package \code{main}:

\begin{salam}
package main

func main:
    println "Hello from package main"
end
\end{salam}

\begin{progout}
Hello from package main
\end{progout}

A file with no \code{package} line belongs to \code{main} by default which is
why the small single-file programs earlier in this book never needed one. As soon
as you split code across files, naming the package makes the structure explicit.

\subsection{Exposing names with \texttt{pub}}

By default, a top-level function, struct, or type is \textbf{private} to its
package: other files cannot see it. To make something visible to code that imports
your module, mark it \code{pub} (short for \emph{public}):

\begin{salam}
// file: geometry.salam
package geometry

pub func area_of_circle(r : f64) : f64:
    ret 3.14159 * r * r
end

func helper() : f64:        // no 'pub' -- private to this file
    ret 0.0
end
\end{salam}

Another file can then import \code{geometry.salam} and use its public functions,
while \code{helper} stays hidden:

\begin{salam}
// file: main.salam
package main
import "geometry.salam"

func main:
    println geometry.area_of_circle(2.0)
end
\end{salam}

This is \textbf{encapsulation}: a module presents a small, deliberate public
surface the things it promises to support while keeping its internal
helpers private and free to change. The same \code{pub} keyword marks the public
parts of the standard library itself; open any \code{std} module and you will see
\code{pub func} and \code{pub struct} on exactly the names you are meant to use.

\begin{tip}
Make things \code{pub} sparingly. The smaller a module's public surface, the
easier it is to use correctly and the freer you are to rewrite its insides without
breaking anyone. Start everything private; promote a name to \code{pub} only when
another file genuinely needs it.
\end{tip}

\section{A tour of the standard library's layout}

It helps to know roughly what is available before you need it. The standard
library is a collection of modules under \code{std}, each in its own folder. A
selection:

\begin{center}
\begin{tabular}{@{}ll@{}}
\toprule
\textbf{Module} & \textbf{Purpose} \\
\midrule
\code{io}, \code{console} & printing and reading input, files \\
\code{fmt}    & formatting values into strings \\
\code{math}   & numeric functions and constants \\
\code{str}    & string operations \\
\code{rand}   & random numbers \\
\code{time}   & dates, timestamps, sleeping \\
\code{collections} & generic containers (\code{Vector}, \code{HashMap}, \dots) \\
\code{os}, \code{fs} & the operating system and file system \\
\code{json}, \code{csv} & data formats \\
\code{regex}  & regular expressions \\
\code{option} & the \code{Option} type \\
\code{mem}    & manual memory management \\
\code{sync}, \code{thread} & concurrency primitives \\
\code{testing} & assertions for unit tests \\
\code{net}, \code{tcp}, \code{http} & networking \\
\bottomrule
\end{tabular}
\end{center}

The next two chapters take a proper tour: Chapter~16 covers the everyday modules
(\code{io}, \code{fmt}, \code{math}, \code{str}, \code{rand}, \code{time}), and
Chapter~17 covers data structures and the world outside your program
(\code{collections}, \code{os}, \code{fs}, \code{json}, and friends).

\begin{deep}[In Depth: qualified names and the generated C]
Qualification is not only for readability it also guarantees that names from
different modules never collide in the compiled program. When \code{math.Sqrt} is
compiled, it becomes a uniquely named C function (something like
\code{\_Salam\_math\_Sqrt\_f64}) that cannot clash with a \code{Sqrt} from any
other module or with your own code. The package system, in other words, gives you
human-friendly short names at the source level and collision-proof unique names
underneath the best of both worlds, with no run-time cost.
\end{deep}

\section{Chapter summary}

\begin{itemize}[leftmargin=1.4em]
  \item \code{import "name"} brings in a standard-library module; call its
        functions qualified, as \code{name.func(...)}.
  \item Import several modules at once by listing their paths after one
        \code{import}; give any import an \code{alias} by writing the alias
        before its path.
  \item \code{import "math"} resolves to \code{std/math/math.salam}; an import
        ending in \code{.salam} is a file relative to the importer (your own
        modules).
  \item Every file has a \code{package}; the entry point is in \code{main} (the
        default when omitted).
  \item Top-level names are private to their package unless marked \code{pub}.
  \item \code{print}, \code{println}, and \code{len} are built in and need no
        import.
\end{itemize}

\section{Exercises}

\exercise Write a program that imports \code{math} and \code{str} on a single
\code{import} line and prints the square root of 2 alongside the word
\code{"root"} in uppercase.

\exercise Import \code{str} under the alias \code{s} and use it to reverse a word.

\exercise Sketch (on paper) a two-file program: a \code{shapes.salam} module with
a \code{pub} function and a private helper, and a \code{main.salam} that imports
it. Mark which names are visible where.

\exercise Open one of the real \code{std} modules in the Salam source and list
five \code{pub} functions it offers. Try calling two of them.

\exercise Explain in a comment why qualification (\code{math.Max}) lets you keep
your own function called \code{max} without conflict.
